\originalsection{第13次习题课}
\sourceinfo{MIT 18.04, Spring 2018, Recitation 13; 题面 \nolinkurl{mit18_04s18_recit13-handout.pdf}, PDF第1页; 解答 \nolinkurl{mit18_04s18_recit13-solutions.pdf}, PDF第1--3页. 课程作者 Jeremy Orloff, 页内署名 Vishesh Jain, 2018, CC BY-NC-SA 4.0}
\begin{exercise}
\textbf{原题 Rec13.1。}\label{ass:rec13-1}
计算 $\mathcal L(\sin(\omega t);s)$, $\omega\in\R$. 官方解答PDF的题面还问: 对哪些 $s$ 可以进行该计算?
\end{exercise}
\begin{solution}
据官方解答翻译补全 (解答PDF第1页); $\omega=0$ 的退化情形为编者补充 (AI). 当 $\Re s>0$ 时, 两个指数变换均绝对收敛, 由线性性
\[
\mathcal L(\sin\omega t;s)=\frac1{2\ii}\left(\frac1{s-\ii\omega}-\frac1{s+\ii\omega}\right)=\frac{\omega}{s^2+\omega^2}.
\]
若 $\omega\ne0$, 通常广义积分的收敛半平面为 $\Re s>0$; 在虚轴上, 两个振荡项的原函数无极限, 在 $s=\pm\ii\omega$ 还出现线性项, 所以有理式的亚纯延拓不能当作原积分的值. 在 $\Re s<0$ 时振荡振幅指数增长, 原函数仍无极限. 若 $\omega=0$, 被变换函数恒为0, 积分对所有 $s\in\C$ 均为0; 尤其 $s=0$ 不能保留未定义形式 $0/0$.
\end{solution}
\begin{exercise}
\textbf{原题 Rec13.2。}\label{ass:rec13-2}
假设 $f$ 有指数类型 $a$. 对 $\Re s>a$, 证明 $\mathcal L(f';s)=s\mathcal L(f;s)-f(0)$. 若 $f'$ 也有指数类型 $a$, 证明 $\mathcal L(f'';s)=s^2\mathcal L(f;s)-sf(0)-f'(0)$.
\end{exercise}
\begin{solution}
据官方解答翻译补全 (解答PDF第1页); 分部积分正则性与收敛含义为编者补充 (AI). 指数类型在这里指存在 $M$ 使 $|f(t)|\le M\ee^{at}$, $t\ge0$. 还须假设可分部积分, 例如第一式 $f\in C^1([0,\infty))$, 第二式 $f\in C^2([0,\infty))$; 局部绝对连续的相应版本也足够. 先在有限区间计算,
\[
\int_0^R f'(t)\ee^{-st}\dd t=f(R)\ee^{-sR}-f(0)+s\int_0^R f(t)\ee^{-st}\dd t.
\]
因 $|f(R)\ee^{-sR}|\le M\ee^{-(\Re s-a)R}\to0$, 右侧存在极限, 得到第一式. 仅由 $f$ 的增长界不一定得到 $f'$ 的绝对可积性, 此时公式按广义积分理解; 若采用只允许绝对收敛的变换定义, 应另假设 $f'$ 有适当指数增长界. 第二式在题给 $f'$ 的指数界下对 $f'$ 重复第一式, 再代入 $\mathcal L(f';s)$, 得 $\mathcal L(f'';s)=-f'(0)+s(s\mathcal L(f;s)-f(0))$. 第二阶变换同样首先按广义积分成立; 若额外有 $f''$ 的指数界, 则也绝对收敛.
\end{solution}
\begin{exercise}
\textbf{原题 Rec13.3。}\label{ass:rec13-3}
假设 $f$ 有指数类型 $a$, $\Re s>a$. 证明 $\mathcal L(tf(t);s)=-\frac{\dd}{\dd s}\mathcal L(f(t);s)$. 据此求所有整数 $n\ge0$ 的 $\mathcal L(t^n;s)$, $\Re s>0$.
\end{exercise}
\begin{solution}
据官方解答翻译补全 (解答PDF第1--2页); 求导与积分交换的估计为编者补充 (AI). 假设 $f$ 局部可积, $|f(t)|\le M\ee^{at}$. 固定 $s_0$ 满足 $\Re s_0>a$, 取 $\delta=(\Re s_0-a)/2>0$. 在 $|s-s_0|<\delta$ 中, $|t f(t)\ee^{-st}|\le Mt\ee^{-\delta t}$, 右侧可积. 对复差商用 $|(\ee^{-ht}-1)/h|\le t\ee^{|h|t}$, 并将 $s$ 与 $h$ 限制在更小邻域, 得到同类可积控制函数. 控制收敛允许将导数移入积分, 故
\[
\frac{\dd}{\dd s}\mathcal L(f;s)=\int_0^\infty -tf(t)\ee^{-st}\dd t=-\mathcal L(tf;s).
\]
从 $\mathcal L(1;s)=1/s$ 开始反复求导, 得 $\mathcal L(t^n;s)=(-1)^n(\dd/\dd s)^n(1/s)=n!/s^{n+1}$. 对每个 $\Re s>0$, 多项式乘指数及其各次 $s$ 导数都可由 $C t^k\ee^{-\delta t}$ 控制, 所以反复交换合法. 无须把 $t^n$ 错认为有统一常数界的指数类型0函数.
\end{solution}
\begin{exercise}
\textbf{原题 Rec13.4。}\label{ass:rec13-4}
解释以下各对函数为何有相同的单侧 Laplace 变换. 4.1. $f(t)=1$ 对所有 $t$ 成立; $u(t)=1$ 当 $t>0$, $u(t)=0$ 当 $t<0$. 4.2. $f(t)=\ee^{at}$ 对所有 $t$ 成立; $g(t)=\ee^{at}$ 当 $t\ne2$, 而 $g(2)=0$.
\end{exercise}
\begin{solution}
据官方解答翻译补全 (解答PDF第2页); 单点取值说明为编者补充 (AI). 4.1单侧变换只积分 $t\ge0$, 负半轴上的差别不参与积分. 原题未定义 $u(0)$, 可以任意指定有限值, 因单点不影响积分. 两个变换在 $\Re s>0$ 均为 $1/s$. 4.2两函数只在 $t=2$ 一个点不同, 被积函数差值也只在这个点非零, 故其积分为0. 当 $a\in\R$ 时, 两者在 $\Re s>a$ 都有变换 $1/(s-a)$; 若允许复参数 $a$, 相应条件为 $\Re s>\Re a$.
\end{solution}
\begin{exercise}
\textbf{原题 Rec13.5。}\label{ass:rec13-5}
用 Laplace 变换和部分分式求解 $x''+8x'+7x=\ee^{-2t}$, 初值 $x(0)=0$, $x'(0)=1$.
\end{exercise}
\begin{solution}
据官方解答翻译补全 (解答PDF第2--3页); 原解答末行误用 $X$ 表示时域函数, 以下改为 $x(t)$, 并补直接核验 (AI). 设 $X(s)=\mathcal L(x;s)$. 对方程取变换并代入初值, 得 $(s^2X-1)+8sX+7X=1/(s+2)$, 所以
\begin{align*}
X(s)&=\frac1{(s+2)(s+7)(s+1)}+\frac1{(s+7)(s+1)}\\
&=-\frac{1}{5(s+2)}+\frac{1}{30(s+7)}+\frac{1}{6(s+1)}
-\frac{1}{6(s+7)}+\frac{1}{6(s+1)}\\
&=\frac{1}{3(s+1)}-\frac{1}{5(s+2)}-\frac{2}{15(s+7)}.
\end{align*}
逐项逆变换得到 $x(t)=\ee^{-t}/3-\ee^{-2t}/5-2\ee^{-7t}/15$. 初值核验为 $x(0)=1/3-1/5-2/15=0$, $x'(0)=-1/3+2/5+14/15=1$. 算子 $(D+1)(D+7)$ 消去 $\ee^{-t}$ 与 $\ee^{-7t}$, 作用在 $-\ee^{-2t}/5$ 上的系数为 $-((-2+1)(-2+7))/5=1$, 正好给 $\ee^{-2t}$. 因而该函数确实满足原方程和初值.
\end{solution}
